\GalSourcePageBoundary{070}{L0069}{41}{41}

vement ceux de l'équation proposée, quand on lui adjoindra successivement
$r,r',r'',\ldots$.

2\textsuperscript{o} Nous avons vu plus haut que toutes les valeurs de $V$ étaient
des fonctions rationnelles les unes des autres. D'après cela, supposons
que, $V$ étant une racine de $f(V,r)=0$, $F(V)$ en soit une
autre; il est clair que de même si $V'$ est une racine de $f(V,r')=0$,
$F(V')$ en sera une autre; car on aura
\[
f[F(V),r]=\text{une fonction divisible par }f(V,r).
\]

Donc (\emph{lemme I})
\[
f[F(V'),r']=\text{une fonction divisible par }f(V',r').
\]

Cela posé, je dis que l'on obtient le groupe relatif à $r'$ en opérant
partout dans le groupe relatif à $r$ une même substitution de
lettres.

En effet, si l'on a, par exemple,
\[
\varphi_{\mu}F(V)=\varphi_{\nu}(V),
\]
on aura encore (\emph{lemme I}),
\[
\varphi_{\mu}F(V')=\varphi_{\nu}(V').
\]
Donc, pour passer de la permutation $[F(V)]$ à la permutation
$[F(V')]$, il faut faire la même substitution que pour passer de la
permutation $(V)$ à la permutation $(V')$.

Le théorème est donc démontré.

\begin{center}{\bfseries PROPOSITION III.}\end{center}
\GalSourceErrorMark{W09-SE003}

\noindent\textsc{Théorème.} --- \emph{Si l'on adjoint à une équation toutes les racines
d'une équation auxiliaire, les groupes dont il est question
dans le théorème II jouiront, de plus, de cette propriété que
les substitutions sont les mêmes dans chaque groupe.}

On trouvera la démonstration\textsuperscript{(*)}.

\GalHistoricalNote{\textsuperscript{(*)} Dans le manuscrit, l'énoncé du théorème qu'on vient de lire se trouve en
marge et en remplace un autre que Galois avait écrit avec sa démonstration sous
le même titre: \emph{Proposition III}. Voici le texte primitif: \textsc{Théorème.} --- Si l'équa-}
